% ===================================================================
% Paper M3 — Which axioms survive: an audit of the twenty axioms of the
% programme against the four.  DRAFT v0, 10 Nov 2026 (architect's draft for
% the PI's read; not canon).
% Compile: xelatex (three passes).  Certificates: MEFProof/Paper_M3/*.lean.
% Premises: M1 and M2 only (island rule).  The earlier statement of the
% twenty axioms is reproduced in Appendix A as data and is not cited.
% ===================================================================
\documentclass[11pt,a4paper]{article}
\usepackage{amsmath,amssymb,amsthm}
\usepackage[margin=25mm]{geometry}
\usepackage{enumitem}
\usepackage{booktabs}
\usepackage{fontspec}
\setmonofont{DejaVu Sans Mono}[Scale=0.78]
\usepackage{listings}
% Listing frames sit inside the text block (Decision Log P15(b)): the frame is
% drawn framesep + rule outside the code, so the code is inset by 5pt each side.
\lstset{basicstyle=\ttfamily\footnotesize,breaklines=true,columns=fullflexible,keepspaces=true,frame=single,framesep=4pt,xleftmargin=5pt,xrightmargin=5pt,extendedchars=true}
\usepackage[hidelinks]{hyperref}
%%=====================================================================
%% mef_macros.tex
%% Canonical macro set for the Master Equation Framework (MEF) series.
%% Version: 1.0 (April 2026)
%%
%% Single source of truth for symbol notation across MEF papers.
%% Derived from main.tex (thesis preamble) plus selected additions
%% from Paper XVI (number sets and rigour-classification marks).
%%
%% Usage: place %%=====================================================================
%% mef_macros.tex
%% Canonical macro set for the Master Equation Framework (MEF) series.
%% Version: 1.0 (April 2026)
%%
%% Single source of truth for symbol notation across MEF papers.
%% Derived from main.tex (thesis preamble) plus selected additions
%% from Paper XVI (number sets and rigour-classification marks).
%%
%% Usage: place %%=====================================================================
%% mef_macros.tex
%% Canonical macro set for the Master Equation Framework (MEF) series.
%% Version: 1.0 (April 2026)
%%
%% Single source of truth for symbol notation across MEF papers.
%% Derived from main.tex (thesis preamble) plus selected additions
%% from Paper XVI (number sets and rigour-classification marks).
%%
%% Usage: place \input{mef_macros.tex} in the preamble of every MEF
%% paper after the package loads but before \begin{document}.  Drift
%% between papers is prevented by editing only this file.
%%
%% Companion file: mef_field_content_preamble.tex (Section 2 reader-
%% orientation block).
%%=====================================================================

%% ---- Quantum Coherence Dilaton and its constituents ----
\newcommand{\QCD}{\Phi_{\mathrm{QC}}}           % Quantum Coherence Dilaton (canonical)
\newcommand{\consciousness}{\omega}             % consciousness coherence field
\newcommand{\zetaval}{12\pi}                    % quantum coherence coupling = chi(K_8) * pi

%% ---- Manifolds ----
\newcommand{\Mtwelve}{\mathcal{M}_{12}}         % 12D parent manifold
\newcommand{\Keight}{K_8}                       % 8D internal manifold
\newcommand{\CPtwo}{\mathbb{CP}^2}              % complex projective plane
\newcommand{\Stwo}{S^2}                         % 2-sphere
\newcommand{\Ttwo}{T^2}                         % 2-torus

%% ---- Discrete groups ----
\newcommand{\Ztwo}{\mathbb{Z}_2}                % pillowcase orbifold group
\newcommand{\Zfour}{\mathbb{Z}_4}

%% ---- Spin^c and bundles ----
\newcommand{\Spinc}{\text{Spin}^c}              % Spin^c structure
\newcommand{\Oline}[1]{\mathcal{O}(#1)}         % line bundle of degree #1

%% ---- Number fields (NB: \Q = mathbb Q is RATIONALS, not the dilaton; \QCD = \mathcal{Q} is the dilaton) ----
\newcommand{\C}{\mathbb{C}}
\newcommand{\R}{\mathbb{R}}
\newcommand{\Z}{\mathbb{Z}}
\newcommand{\Q}{\mathbb{Q}}

%% ---- Effective couplings, scales, moduli ----
\newcommand{\Veff}{V_{\text{eff}}}
\newcommand{\VCKM}{V_{\text{CKM}}}
\newcommand{\lambdaC}{\lambda_C}
\newcommand{\tautwo}{\tau_2}                    % stabilised aspect modulus
\newcommand{\thetaW}{\theta_W}                  % Weinberg angle
\newcommand{\Rw}{R_w}                           % S^2 radius (warped fibre)
\newcommand{\Rc}{R_c}                           % CP^2 radius (Calabi-like)
\newcommand{\Rfibre}{\mathcal{R}_{\text{fibre}}}
\newcommand{\Cfibre}{C_{\mathrm{fibre}}}
\newcommand{\Rsix}{\mathcal{R}_6}
\newcommand{\Fsq}{\mathcal{F}^2}
\newcommand{\centcorr}{\mathcal{C}_{32}}

%% ---- Spin^c holonomy phases ----
\newcommand{\Phiomega}{\phi_\consciousness}     % omega-axis holonomy
\newcommand{\Phipsi}{\phi_\psi}                 % psi-axis holonomy

%% ---- Arithmetic ----
\newcommand{\dodd}{d_{\mathrm{odd}}}            % odd divisor function

%% ---- Operators ----
\newcommand{\Tr}{\mathrm{Tr}}
\newcommand{\Vol}{\mathrm{Vol}}
\newcommand{\vol}{\mathrm{vol}}

%% ---- Rigour-classification marks ----
\newcommand{\rigmark}{\hfill\textbf{[Rigorous]}}
\newcommand{\dermark}{\hfill\textbf{[Derived]}}
\newcommand{\motmark}{\hfill\textbf{[Motivated]}}
\newcommand{\gapmark}{\hfill\textbf{[Gap]}}
\newcommand{\rigder}{\hfill\textbf{[Rigorous / Derived]}}
\newcommand{\dergap}{\hfill\textbf{[Derived / Gap]}}

%% ---- End of mef_macros.tex ----
 in the preamble of every MEF
%% paper after the package loads but before \begin{document}.  Drift
%% between papers is prevented by editing only this file.
%%
%% Companion file: mef_field_content_preamble.tex (Section 2 reader-
%% orientation block).
%%=====================================================================

%% ---- Quantum Coherence Dilaton and its constituents ----
\newcommand{\QCD}{\Phi_{\mathrm{QC}}}           % Quantum Coherence Dilaton (canonical)
\newcommand{\consciousness}{\omega}             % consciousness coherence field
\newcommand{\zetaval}{12\pi}                    % quantum coherence coupling = chi(K_8) * pi

%% ---- Manifolds ----
\newcommand{\Mtwelve}{\mathcal{M}_{12}}         % 12D parent manifold
\newcommand{\Keight}{K_8}                       % 8D internal manifold
\newcommand{\CPtwo}{\mathbb{CP}^2}              % complex projective plane
\newcommand{\Stwo}{S^2}                         % 2-sphere
\newcommand{\Ttwo}{T^2}                         % 2-torus

%% ---- Discrete groups ----
\newcommand{\Ztwo}{\mathbb{Z}_2}                % pillowcase orbifold group
\newcommand{\Zfour}{\mathbb{Z}_4}

%% ---- Spin^c and bundles ----
\newcommand{\Spinc}{\text{Spin}^c}              % Spin^c structure
\newcommand{\Oline}[1]{\mathcal{O}(#1)}         % line bundle of degree #1

%% ---- Number fields (NB: \Q = mathbb Q is RATIONALS, not the dilaton; \QCD = \mathcal{Q} is the dilaton) ----
\newcommand{\C}{\mathbb{C}}
\newcommand{\R}{\mathbb{R}}
\newcommand{\Z}{\mathbb{Z}}
\newcommand{\Q}{\mathbb{Q}}

%% ---- Effective couplings, scales, moduli ----
\newcommand{\Veff}{V_{\text{eff}}}
\newcommand{\VCKM}{V_{\text{CKM}}}
\newcommand{\lambdaC}{\lambda_C}
\newcommand{\tautwo}{\tau_2}                    % stabilised aspect modulus
\newcommand{\thetaW}{\theta_W}                  % Weinberg angle
\newcommand{\Rw}{R_w}                           % S^2 radius (warped fibre)
\newcommand{\Rc}{R_c}                           % CP^2 radius (Calabi-like)
\newcommand{\Rfibre}{\mathcal{R}_{\text{fibre}}}
\newcommand{\Cfibre}{C_{\mathrm{fibre}}}
\newcommand{\Rsix}{\mathcal{R}_6}
\newcommand{\Fsq}{\mathcal{F}^2}
\newcommand{\centcorr}{\mathcal{C}_{32}}

%% ---- Spin^c holonomy phases ----
\newcommand{\Phiomega}{\phi_\consciousness}     % omega-axis holonomy
\newcommand{\Phipsi}{\phi_\psi}                 % psi-axis holonomy

%% ---- Arithmetic ----
\newcommand{\dodd}{d_{\mathrm{odd}}}            % odd divisor function

%% ---- Operators ----
\newcommand{\Tr}{\mathrm{Tr}}
\newcommand{\Vol}{\mathrm{Vol}}
\newcommand{\vol}{\mathrm{vol}}

%% ---- Rigour-classification marks ----
\newcommand{\rigmark}{\hfill\textbf{[Rigorous]}}
\newcommand{\dermark}{\hfill\textbf{[Derived]}}
\newcommand{\motmark}{\hfill\textbf{[Motivated]}}
\newcommand{\gapmark}{\hfill\textbf{[Gap]}}
\newcommand{\rigder}{\hfill\textbf{[Rigorous / Derived]}}
\newcommand{\dergap}{\hfill\textbf{[Derived / Gap]}}

%% ---- End of mef_macros.tex ----
 in the preamble of every MEF
%% paper after the package loads but before \begin{document}.  Drift
%% between papers is prevented by editing only this file.
%%
%% Companion file: mef_field_content_preamble.tex (Section 2 reader-
%% orientation block).
%%=====================================================================

%% ---- Quantum Coherence Dilaton and its constituents ----
\newcommand{\QCD}{\Phi_{\mathrm{QC}}}           % Quantum Coherence Dilaton (canonical)
\newcommand{\consciousness}{\omega}             % consciousness coherence field
\newcommand{\zetaval}{12\pi}                    % quantum coherence coupling = chi(K_8) * pi

%% ---- Manifolds ----
\newcommand{\Mtwelve}{\mathcal{M}_{12}}         % 12D parent manifold
\newcommand{\Keight}{K_8}                       % 8D internal manifold
\newcommand{\CPtwo}{\mathbb{CP}^2}              % complex projective plane
\newcommand{\Stwo}{S^2}                         % 2-sphere
\newcommand{\Ttwo}{T^2}                         % 2-torus

%% ---- Discrete groups ----
\newcommand{\Ztwo}{\mathbb{Z}_2}                % pillowcase orbifold group
\newcommand{\Zfour}{\mathbb{Z}_4}

%% ---- Spin^c and bundles ----
\newcommand{\Spinc}{\text{Spin}^c}              % Spin^c structure
\newcommand{\Oline}[1]{\mathcal{O}(#1)}         % line bundle of degree #1

%% ---- Number fields (NB: \Q = mathbb Q is RATIONALS, not the dilaton; \QCD = \mathcal{Q} is the dilaton) ----
\newcommand{\C}{\mathbb{C}}
\newcommand{\R}{\mathbb{R}}
\newcommand{\Z}{\mathbb{Z}}
\newcommand{\Q}{\mathbb{Q}}

%% ---- Effective couplings, scales, moduli ----
\newcommand{\Veff}{V_{\text{eff}}}
\newcommand{\VCKM}{V_{\text{CKM}}}
\newcommand{\lambdaC}{\lambda_C}
\newcommand{\tautwo}{\tau_2}                    % stabilised aspect modulus
\newcommand{\thetaW}{\theta_W}                  % Weinberg angle
\newcommand{\Rw}{R_w}                           % S^2 radius (warped fibre)
\newcommand{\Rc}{R_c}                           % CP^2 radius (Calabi-like)
\newcommand{\Rfibre}{\mathcal{R}_{\text{fibre}}}
\newcommand{\Cfibre}{C_{\mathrm{fibre}}}
\newcommand{\Rsix}{\mathcal{R}_6}
\newcommand{\Fsq}{\mathcal{F}^2}
\newcommand{\centcorr}{\mathcal{C}_{32}}

%% ---- Spin^c holonomy phases ----
\newcommand{\Phiomega}{\phi_\consciousness}     % omega-axis holonomy
\newcommand{\Phipsi}{\phi_\psi}                 % psi-axis holonomy

%% ---- Arithmetic ----
\newcommand{\dodd}{d_{\mathrm{odd}}}            % odd divisor function

%% ---- Operators ----
\newcommand{\Tr}{\mathrm{Tr}}
\newcommand{\Vol}{\mathrm{Vol}}
\newcommand{\vol}{\mathrm{vol}}

%% ---- Rigour-classification marks ----
\newcommand{\rigmark}{\hfill\textbf{[Rigorous]}}
\newcommand{\dermark}{\hfill\textbf{[Derived]}}
\newcommand{\motmark}{\hfill\textbf{[Motivated]}}
\newcommand{\gapmark}{\hfill\textbf{[Gap]}}
\newcommand{\rigder}{\hfill\textbf{[Rigorous / Derived]}}
\newcommand{\dergap}{\hfill\textbf{[Derived / Gap]}}

%% ---- End of mef_macros.tex ----


\theoremstyle{plain}
\newtheorem{theorem}{Theorem}[section]
\newtheorem{proposition}[theorem]{Proposition}
\newtheorem{lemma}[theorem]{Lemma}
\newtheorem{corollary}[theorem]{Corollary}
\theoremstyle{definition}
\newtheorem{definition}[theorem]{Definition}
\newtheorem{axiom}{Axiom}
\theoremstyle{remark}
\newtheorem{remark}[theorem]{Remark}

\newcommand{\Hquat}{\mathbb{H}}
\newcommand{\rigR}{[\textup{R}]}
\newcommand{\rigD}{[\textup{D}]}
\newcommand{\rigM}{[\textup{M}]}
\newcommand{\rigG}{[\textup{G}]}
\newcommand{\Max}[1]{$\mathbb{M}$-#1}
\newcommand{\Sig}[1]{$\Sigma$-#1}

\title{Which axioms survive: an audit of the twenty axioms of the programme against the four}
\author{Dhiren J. Master}
\date{10 November 2026 (draft v0)}

\begin{document}
\maketitle

\begin{abstract}
The earlier work of the Master Equation programme was stated on twenty axioms: twelve describing a twelve-dimensional geometry and its fields, and eight describing the mathematics by which results were obtained from it. The first two papers of this series derived that geometry and its field equation from four sentences about what is observed, together with a principle of minimality. We audit the twenty against those two papers and classify each as a theorem already proved, a target owed by a named later paper, an axiom, or not an axiom. Eight are theorems of the first two papers, with corrections where the earlier wording does not survive. Seven are targets, and for each the paper that owes it is named. Three are not axioms: a method, a statement about provability, and an interpretation. Two are one statement, and it is the only one that the four axioms do not reach: that what is observed is a shadow, which does not carry all the information of the object that casts it. We state it as a fifth axiom, as an existence statement without uniqueness. The series therefore rests on five sentences, one principle and one measured scale.
\end{abstract}

\section{What is observed, and what this paper does}\label{sec:intro}

The purpose of this paper is to say how many axioms the series needs. The first paper, M1, began from four observations that need no instrument, took them as axioms, and derived a twelve-dimensional geometry from them. The second, M2, derived the field equation on that geometry from the same four. The programme these papers rebuild was not written that way. It was stated on twenty axioms, and this paper goes through the twenty one at a time and asks of each whether it is still needed.

\subsection{How the twenty arose}\label{sec:history}

The twenty were not chosen at the outset. In the earlier work more than ninety quantitative statements were obtained, across particle physics, cosmology and number theory, and certain patterns kept recurring in the derivations: a fold with four corners, a drive that stretches and squeezes, a field and its mirror image, a singular object that became regular when seen in more dimensions. Each recurring pattern was written down as an axiom, and the list grew to twenty. A list assembled in that way records what was found useful; it does not show which of its entries are independent. The count of statements obtained is given here as history and is not a result of this paper.

The first two papers of the series change the position. If the geometry is derived, every axiom that describes the geometry is a theorem or it is false; if the field equation is derived, every axiom that describes the fields is a theorem or it is false. What remains after that subtraction is the set of statements the four axioms do not reach, and the purpose of the audit is to find that set.

\subsection{The four axioms}\label{sec:four}

We restate the four as M1 did.

\begin{axiom}[closure]\label{ax:A1}
The universe closes on itself, so that it can be aware of itself.
\end{axiom}

\begin{axiom}[growth and decay]\label{ax:A2}
The universe grows and decays as a function of itself.
\end{axiom}

\begin{axiom}[order]\label{ax:A3}
The order of events matters, so that each instance is unique and cannot be foretold.
\end{axiom}

\begin{axiom}[drive]\label{ax:A4}
The universe is in constant change; it never stands still.
\end{axiom}

One remark connects the four, and it is a theorem of M2 rather than a decoration. The operation that exchanges growth with decay is a reversal of the drive together with a quarter-turn of the internal fibre (M2 Proposition~2.6). A quarter-turn done twice is a half-turn, and in the complex coordinate of the plane a half-turn is multiplication by $e^{i\pi}=-1$. So two exchanges of growth and decay are one fold of the fibre, and the fold is what makes the pillowcase of M1: the exponential of the second axiom, the turn of the fourth and the angle $\pi$ of the corners meet in the one relation $S^{2}=-I$ (certificate \texttt{L4\_S\_sq\_fold} of M2; re-derived as \texttt{L3\_exchange\_squared}). \rigR

\subsection{The result in plain words}\label{sec:plain}

Of the twenty, eight are already proved in M1 and M2. Seven are statements about spectra, partition functions and the passage between dimensions that later papers of the series must prove, and we name the paper that owes each. Three are not axioms at all. The remaining two are the same statement written twice, once in words and once in mathematics, and it is the one statement the four axioms do not supply.

That statement can be put as plainly as the four. We can see the shadow, but the shadow is missing information about what it is a shadow of. Two shadows can be identical and be cast by different objects; one object can cast shadows that look nothing alike. The four axioms say what the object is. They do not say that everything we observe, including what looks broken or singular, is a shadow of that object. That is the fifth axiom.

\subsection{What this paper proves, and how it is organised}\label{sec:results}

The results are these. The twenty axioms fall into four classes (Definition~\ref{def:classes}), and the class of each is given in Theorem~\ref{thm:classification} with its reason. The fifth axiom is stated in words (Axiom~\ref{ax:A5}) and formally (Definition~\ref{def:A5}), as the existence of a regular twelve-dimensional configuration above every observed one; that the passage down loses information is a remark with its proof (Remark~\ref{rem:loss}), and uniqueness is not asserted. The smallest case, the corner of the pillowcase, is worked by hand (Proposition~\ref{prop:smallest}). The minimal set of the series is stated in one sentence (Theorem~\ref{thm:minimal}).

The paper is organised as follows. Section~\ref{sec:setting} recalls what M1 and M2 established, in the form the audit uses, and defines the classes. Section~\ref{sec:theorems} treats the axioms that are theorems, Section~\ref{sec:targets} those that are targets, Section~\ref{sec:axiom} the fifth axiom, and Section~\ref{sec:notaxioms} those that are not axioms. Section~\ref{sec:result} states the minimal set and the list of targets. Appendix~\ref{app:axioms} reproduces the twenty axioms as they were stated; Appendix~\ref{app:labels} is the cumulative table of symbols; Appendix~\ref{app:lean} contains the certificates. As in M1 and M2, this introduction is intended to be followed by a first-year student; from Section~\ref{sec:setting} the register is that of a formal paper. Rigour is tagged \rigR\ (proved, or a cited classical theorem), \rigD\ (derived under a stated reading or under minimality), \rigM\ (motivated) and \rigG\ (open), as in M1.

The deflationary reading is stated here, where it belongs. This paper proves nothing new about the geometry or the field equation. It is bookkeeping: it shows which of the earlier axioms are consequences of results already proved, and it identifies what is left. A verdict of ``theorem'' below is only as strong as the theorem of M1 or M2 it cites, with the readings and the minimal choices that theorem carries. A verdict of ``target'' is a statement that something has not yet been proved in the series.

\section{The setting: what M1 and M2 established, and the classes}\label{sec:setting}

\subsection{The results the audit uses}\label{sec:recall}

M1 formalises the four axioms as conditions on a compact Riemannian geometry $(T\times\Sigma)\times F$ under four numbered readings (M1 Definitions~2.1, 2.2(i), 2.5, 2.6), defines minimality by four criteria applied in order (M1 Definition~3.2), and proves that the minimal geometry is unique (M1 Theorem~8.1):
\begin{equation}\label{eq:M12}
\Mtwelve=(S^{1}_{\mathrm{Sol}}\times S^{3})\times\Keight,\qquad
\Keight=\bigl[(\CPtwo\times\Stwo)\times_{w}(\Ttwo/\Ztwo)\bigr]^{\Spinc},
\end{equation}
with monodromy $M_{0}=\bigl(\begin{smallmatrix}2&1\\1&1\end{smallmatrix}\bigr)$, determinant line $\Oline{3}$, $\chi(\Keight)=12$, dimension twelve and no continuous modulus, the overall scale being a declared input. The steps used below are: the drive is unique at the least entropy, with eigenvalues $\varphi^{\pm2}$ and entropy $2\log\varphi$ (M1 Theorem~4.2), and fixes one corner while cycling the other three (M1 Proposition~4.4); the fibre algebra is the quaternions (M1 Proposition~5.2); the fibre is a torus folded by $z\mapsto-z$ to the pillowcase, a sphere with four cone points of angle $\pi$ (M1 Proposition~5.5, Theorem~5.7), carrying $\R$, $\mathbb{C}_{i}$, $\mathbb{C}_{j}$, $\mathbb{C}_{k}$ at the corners and $\Hquat$ elsewhere (M1 Proposition~5.6); the dimension is twelve (M1 Theorem~6.1); the base is $\CPtwo\times\Stwo$ with a $\Spinc$ structure forced by $w_{2}(\CPtwo)\neq0$ (M1 Propositions~7.1 to~7.3) and corner lift signs $(-1)^{mn}$ (M1 Corollary~7.4); the spatial factor is $S^{3}$ (M1 Proposition~7.5); the aspect ratio of the fibre is $\tautwo=1.8466$ (M1 Proposition~8.2, whose stationarity conditions are hypotheses owed by a later paper).

M2 adds a fifth reading, that the metric is positive-definite and time is the orbit parameter of the drive (M2 Definition~2.1), and states the standing principle that wherever the axioms leave a choice the minimal option is taken and named. It proves that the two sheets of the growth law are one field $\psi$, a section of the Clifford module $M\cong\Hquat^{32}$, with shadow $\bar\psi=\hat\tau\psi$ (M2 Definition~2.2); that $\omega=\psi+\bar\psi$ and $\QCD=\omega\lvert\psi-\bar\psi\rvert^{2}$ is the unique lowest-order even scalar vanishing on both degenerate loci (M2 Definition~2.3, Proposition~2.4); that the exchange $\tau$ is the quarter-turn flow reversal with $\tau^{2}$ the fold (M2 Proposition~2.6) and that $\bar\psi$ has the opposite handedness to $\psi$ (M2 Theorem~2.7); that the field equation is
\begin{equation}\label{eq:master}
G_{MN}=8\pi G_{\mathrm{eff}}\,T_{MN}+\zeta\,T^{\QCD}_{MN}+\Lambda_{\mathrm{eff}}\,g_{MN}
\end{equation}
(M2 Theorem~3.1); that $\zeta=\chi(\Keight)\pi=12\pi$ (M2 Theorem~4.1, Corollary~4.2); and that $\omega$ vanishes at the locked phase and averages to zero at the uncorrelated phase (M2 Propositions~5.1, 5.3).

\subsection{The twenty axioms}\label{sec:twenty}

The earlier statement has two tiers. The first has twelve entries, \Sig{0} to \Sig{8} with \Sig{P}, \Sig{S} and \Sig{T}; \Sig{2} has eight parts, a to h, one for each component of the geometry, and \Sig{3} has a part 3a. The second tier has eight, \Max{1} to \Max{8}. They are reproduced in Appendix~\ref{app:axioms} with their source lines; the count is certified in \texttt{L2\_twenty}. Most entries consist of a headline statement followed by consequences for physics that the headline is said to produce. We call these consequences \emph{riders}.

\begin{definition}[the classes]\label{def:classes}
An axiom of the earlier statement is of class \textbf{T} if its headline is a theorem of M1 or M2, the theorem being named; of class \textbf{TT} if it is a theorem target, a statement not proved in M1 or M2 and owed by a named later paper of the series; of class \textbf{A} if it is not reached by the four axioms and is declared; and of class \textbf{N} if it is not an axiom: a method, a statement about the theory, or an interpretation. Each axiom has exactly one class. The class is that of the headline; riders are listed with their own status. Where an entry is a single block with no headline, it is of class T only if every assertive clause of it is a theorem of M1 or M2.
\end{definition}

The rule in the last sentence is the minimal one: it never credits an entry with more than has been proved. \rigD\ (minimality)

\begin{theorem}[the classification]\label{thm:classification}
Under Definition~\ref{def:classes} the twenty axioms are classified as in Table~\ref{tab:classes}: eight of class T, seven of class TT, two of class A which are one statement, and three of class N.
\end{theorem}

\begin{table}[t]
\begin{center}\small
\caption{The classification of the twenty axioms.}\label{tab:classes}
\begin{tabular}{@{}p{0.27\textwidth}lp{0.58\textwidth}@{}}
\toprule
Axiom & Class & Reason, and where it is argued \\
\midrule
\Sig{0} Universal Symmetry & T & M2 Proposition~2.6, Theorem~2.7; exact for the exponent, not for the periods (Section~\ref{sec:sigma0}) \\
\Sig{1} The Manifold & T & M1 Theorem~8.1 (Section~\ref{sec:shapes}) \\
\Sig{2} The Fundamental Shapes & T & M1 Theorems~4.2, 5.7, 7.6, 8.1 and Propositions~7.1 to~7.5, part by part (Section~\ref{sec:shapes}) \\
\Sig{3} The Pillowcase, with 3a & T & in part: M1 Theorem~5.7, Proposition~5.6; M2 Definition~2.2 (Section~\ref{sec:sigma3}) \\
\Sig{4} The Three Fields & T & M2 Definitions~2.2, 2.3, Proposition~2.4; one field, not three (Section~\ref{sec:sigma4}) \\
\Sig{5} The Shadow Principle & A & the fifth axiom (Section~\ref{sec:axiom}) \\
\Sig{6} Shape Determinism & T & M1 Definition~3.2(1), Theorem~8.1; one declared scale (Section~\ref{sec:sigma6}) \\
\Sig{7} Dimensional Simplicity & N & an interpretation (Section~\ref{sec:notaxioms}) \\
\Sig{8} The $\omega$-Drive & TT & field-level zeros in M2; spectral statement owed by M4 \\
\Sig{P} Dimensional Permeability & TT & owed by M8 and M11 \\
\Sig{S} Shadow Influence & T & headline: M2 Definition~2.2, Theorem~2.7; channels owed by M5, M17 \\
\Sig{T} Dual Transparency & TT & owed by M8 \\
\Max{1} Observables as spectral data & TT & owed by M4 \\
\Max{2} The Shadow Completion & A & the same statement as \Sig{5} \\
\Max{3} The proof template & N & a method \\
\Max{4} The G\"odel boundary & N & a statement about provability \\
\Max{5} The partition function & TT & owed by M4 and M5 \\
\Max{6} The spectral data & TT & owed by M4 \\
\Max{7} The Sol monodromy & T & M1 Theorem~4.2, with one correction (Section~\ref{sec:m7}) \\
\Max{8} The $\Spinc$ twist & TT & first clause M1 Proposition~7.3; the rest owed by M4, M5 \\
\bottomrule
\end{tabular}
\end{center}
\end{table}

\begin{proof}
The verdicts are argued one by one in Sections~\ref{sec:theorems} to~\ref{sec:notaxioms}. The counts, and the fact that the classes partition the twenty, are \texttt{L2\_class\_counts} and \texttt{L2\_partition}. \rigD
\end{proof}

\section{The axioms that are theorems}\label{sec:theorems}

Each verdict in this section cites a theorem of M1 or M2 and inherits its tag. None is unconditional: each rests on the readings and the minimal choices of the theorem cited, and where the earlier wording says more than the theorem, the excess is named and moved to Section~\ref{sec:targets}.

\subsection{The manifold and its shapes: \texorpdfstring{$\Sigma$}{Sigma}-1 and \texorpdfstring{$\Sigma$}{Sigma}-2}\label{sec:shapes}

\Sig{1} states that the universe is twelve-dimensional with the geometry~\eqref{eq:M12}. This is M1 Theorem~8.1, at \rigD. Two clauses of its wording are not part of the theorem. That the additional dimensions are ``vast, not small'' is a statement about the overall scale, which M1 does not derive and declares as its one input; it is neither a theorem nor a further axiom. That every singular phenomenon is regular in the full geometry is \Sig{5}, treated in Section~\ref{sec:axiom}.

\Sig{2} states that the geometry consists of specific components, each necessary, and lists eight. The headline is M1 Theorem~8.1, whose proof names the axiom and the criterion that force each factor. The parts are as follows, each at the tag of the result cited. The plane $\CPtwo$ is M1 Proposition~7.2, the least direct step of M1 and tagged so there. The sphere $\Stwo$ is M1 Proposition~7.1. The torus is M1 Proposition~5.5, and its aspect ratio is fixed with no modulus by M1 Proposition~8.2; the value is $1.8466$, where the earlier statement has $1.86$. The involution is M1 Proposition~5.5 and Theorem~5.7. The warped product is the second axiom itself under M1's second reading, and its exponent is M2 Theorem~4.1. The $\Spinc$ structure is M1 Proposition~7.3. The time circle with its monodromy is M1 Theorem~4.2. The spatial sphere is M1 Proposition~7.5. The Euler characteristic $3\cdot2\cdot2=12$ is M1 Theorem~7.6.

One refinement of M2 must be carried into every later use of these axioms. The earlier statement speaks of ``the $\Ztwo$ involution'' for two operations. The fold $z\mapsto-z$ makes the pillowcase. What separates the field from its shadow is the exchange $\tau$, a quarter-turn with a reversal of the drive, whose square is the fold (M2 Proposition~2.6). On the four corners the fold fixes every corner, the drive fixes $P_{1}$ and cycles the other three, and the exchange fixes $P_{1}$ and $P_{4}$ and swaps $P_{2}$ with $P_{3}$ (\texttt{L3\_fold\_fixes\_corners}, \texttt{L3\_drive\_corners}, \texttt{L3\_exchange\_corners}).

The riders of \Sig{2} identify the components with physics: the isometries of $\CPtwo$ with the strong force, three generations, the weak mixing angle, the Higgs mechanism, the hierarchy, and others. M1 names no gauge group and no particle, and uses Witten's theorem only as a cross-check. These identifications are targets and are listed in Section~\ref{sec:targets}. One clause of the headline, that whether other geometries could produce other universes is undecidable from within, is a statement about the theory and is of the kind treated in Section~\ref{sec:notaxioms}; M1 itself does not claim that a non-minimal geometry fails the axioms.

\subsection{The pillowcase: \texorpdfstring{$\Sigma$}{Sigma}-3 and 3a}\label{sec:sigma3}

Part 3a states that the fibre algebra is non-commutative over a commutative base and reduces to the real numbers at exactly one of the four corners. This is M1 Proposition~5.6 with Proposition~4.4, at \rigD: the algebra is $\R$ at $P_{1}$, a complex line at each of the other three, and $\Hquat$ at every point with trivial stabiliser.

The headline of \Sig{3} states that the base has two walls and an axis of $\omega$, making four sectors. Part of this is in hand: the pillowcase and its four corners (M1 Theorem~5.7), and two sheets $\psi$ and $\bar\psi$ of opposite handedness (M2 Definition~2.2, Theorem~2.7). The rest is not. In M2, $\omega$ is a field, the even part of the sheet field, and not a coordinate axis; the exchange $\tau$ swaps the corners at the $\psi$ and $\omega$ half-periods; and the word ``wall'' is used without a definition. The two-by-two table of sectors is therefore not a structure that M1 or M2 exhibits. We class \Sig{3} as T for what is proved and record the re-expression of its wall picture in terms of $\tau$ as a target without an owner (Table~\ref{tab:unowned}). \rigD

\subsection{The fields: \texorpdfstring{$\Sigma$}{Sigma}-4}\label{sec:sigma4}

\Sig{4} states that there are three fundamental fields, $\omega$, $\psi$ and $\bar\psi$, coupled through $\QCD$. The fields and the coupling are M2 Definitions~2.2 and~2.3 and Proposition~2.4, with two corrections that M2 settles. There is one field and not three: the shadow and $\omega$ are both determined by $\psi$, and the independent parts are its even and odd parts. And the square in the coupling is the modulus squared, $\QCD=\omega\lvert\psi-\bar\psi\rvert^{2}$, since $\omega$ is quaternion-valued.

One rider does not hold as written. The earlier statement asserts that $\QCD$ is non-zero for a topological reason. In M2 the coupling vanishes at the real corner, wherever the sheets are locked (M2 Proposition~5.1), and wherever they coincide; read pointwise the rider is false. What M2 proves is weaker: a non-zero sheet field of definite handedness cannot equal its shadow (M2 Theorem~2.7). The rider needs a precise statement and a proof where it is first used, and is listed in Table~\ref{tab:unowned}.

\subsection{Symmetry: \texorpdfstring{$\Sigma$}{Sigma}-0}\label{sec:sigma0}

\Sig{0} states that the universe is exactly symmetric across the involution and that every asymmetry is an artefact of observation from one side. The exchange $\tau$ exists and is a symmetry of the Sol block (M2 Proposition~2.6), it is the involution of the second axiom (\rigD, a minimal choice of M2), and the two sectors it exchanges have opposite handedness (M2 Theorem~2.7), so that their handedness sums to zero (\texttt{L3\_chirality\_sum}).

The symmetry is exact for the warp exponent and not for the period lengths. M2 shows that the exponent $12\pi$ is the same for either sheet, and that the aspect ratio $\tautwo$ is not a convention: the decay rate differs by $\tautwo$ between the sheets. So the two sides differ in a real quantity, and ``every asymmetry is an artefact'' does not hold as stated. We class \Sig{0} as T with this caveat, at \rigD. Its riders, that total baryon number and total charge vanish across both sides, are not touched by M1 or M2 (Table~\ref{tab:unowned}).

\subsection{No free parameters: \texorpdfstring{$\Sigma$}{Sigma}-6}\label{sec:sigma6}

\Sig{6} states that everything is fixed by the shapes and that there are no free parameters. Its content is the first criterion of minimality, that a minimal geometry has no modulus (M1 Definition~3.2), together with the theorem that the minimal geometry exists and has none (M1 Theorem~8.1, Proposition~8.2). It is therefore a theorem relative to the four axioms with minimality, at \rigD, and not a further axiom. Two qualifications are part of the verdict. There is one dimensionful input, the overall scale, which M1 declares (\texttt{L2\_inputs}). And the stationarity conditions behind the aspect ratio are hypotheses of M1 owed by M7. The rider, that each named physical quantity is an invariant of the geometry, is a target paper by paper.

\subsection{The drive: \texorpdfstring{$\mathbb{M}$}{M}-7}\label{sec:m7}

\Max{7} states the monodromy, its eigenvalues, its Lyapunov exponent, the fixed point of its M\"obius action, and that it is Anosov. All are M1 Theorem~4.2 and Proposition~4.1, at \rigR\ given the second criterion of minimality, with one correction: the eigenvalues are $\varphi^{2}$ and $\varphi^{-2}$, so the exponent is $2\log\varphi$, as M1 states, and not $\log\varphi$ as the earlier axiom has it. The mixing property and the arrow of time are M2 Proposition~5.3 and Theorem~2.7.

\subsection{The shadow sector: \texorpdfstring{$\Sigma$}{Sigma}-S}\label{sec:sigmaS}

The headline of \Sig{S} states that the second sector is not in a different dimension but is the mirror image of the first within the same geometry. This is M2 Definition~2.2 and Theorem~2.7: the shadow is a section of the same module over the same $\Mtwelve$, with the opposite handedness. Of its three channels of influence, the second, through $\omega$, holds by the definition of $\omega$; the first, through gravity, has its form in~\eqref{eq:master}, where the shadow enters the curvature through $\zeta T^{\QCD}$, and its identification with dark matter is owed by M17; the third, through a completion of a modular object, is owed by M5.

\section{The axioms that are targets}\label{sec:targets}

A target is a statement that M1 and M2 do not prove and that a later paper of the series owes. Table~\ref{tab:targets} names the owing paper for each of the seven axioms of class TT and for the riders of the axioms of class T. The assignment follows the order of the series, in which nothing is used before the paper that derives it.

\begin{table}[t]
\begin{center}\small
\caption{The targets and the papers that owe them.}\label{tab:targets}
\begin{tabular}{@{}p{0.47\textwidth}p{0.12\textwidth}p{0.33\textwidth}@{}}
\toprule
Statement & Owed by & In hand from M1 and M2 \\
\midrule
\Max{6}: the spectrum and its degeneracies $\dodd(n)$ & M4 & nothing \\
\Max{5}: the partition function and its modular type; its completion & M4, M5 & nothing \\
\Max{1}: observables as internal spectral data & M4 & nothing \\
\Max{8}: what the $\Spinc$ twist does to the partition function & M4, M5 & that the twist is forced (M1 Proposition~7.3) \\
\Sig{8}: $\omega$ as a functional of a spectrum & M4 & the two zeros of $\omega$ and one hump (M2 Propositions~5.1, 5.3) \\
\Sig{P}: gauge fields confined; punctures & M8, M11 & the metric and $\QCD$ are fields on $\Mtwelve$ (M2 Definition~2.5) \\
\Sig{T}: exactly two fields cross every boundary & M8 & the two geometric sources in~\eqref{eq:master} \\
\Sig{S}: the shadow as dark matter; the completion & M17, M5 & the headline (Section~\ref{sec:sigmaS}) \\
Riders of \Sig{2}: gauge groups, generations, mixing & M8, M9, M13 & the shapes (Section~\ref{sec:shapes}) \\
Riders of \Sig{2}: hierarchy, vacuum energy & M10 & $\zeta=12\pi$ (M2 Corollary~4.2) \\
Riders of \Sig{2}: fine-structure constant; Higgs & M14, M15 & nothing \\
Stationarity conditions behind $\tautwo$ & M7 & the cubic (M1 Proposition~8.2) \\
Completeness of the geometry; the completeness axiom of the black-hole paper & M11 & nothing stated here \\
\bottomrule
\end{tabular}
\end{center}
\end{table}

Three remarks on the table. First, \Sig{8} is one-sided as it was written: it describes low $\omega$ as ordered and high $\omega$ as interconnected and omits the second zero, at the uncorrelated phase, which M2 derives. Any later statement of it must carry both zeros.

Second, the function $\dodd$ of \Max{6} is the number of odd divisors of $n$, and not their sum as the earlier axiom says. The axiom itself fixes this, since it gives the Dirichlet series as $\zeta(s)^{2}(1-2^{-s})$, and
\[
\sum_{n\ge1}\frac{\dodd(n)}{n^{s}}
=\Bigl(\sum_{m\ge1}\frac1{m^{s}}\Bigr)\Bigl(\sum_{d\ \mathrm{odd}}\frac1{d^{s}}\Bigr)
=\zeta(s)\cdot\zeta(s)\bigl(1-2^{-s}\bigr)
\]
when $\dodd$ counts the odd divisors, each $n$ being written as $md$ with $d$ odd in $\dodd(n)$ ways. The smallest case by hand: $6$ has the odd divisors $1$ and $3$, so $\dodd(6)=2$. \rigR

Third, the completeness of the geometry is not stated in this paper. The question whether the completeness assumed in the treatment of the black-hole interior is a consequence of closure is left to the paper that treats it.

\begin{table}[t]
\begin{center}\small
\caption{Riders with no owing paper yet named in the series.}\label{tab:unowned}
\begin{tabular}{@{}p{0.62\textwidth}p{0.30\textwidth}@{}}
\toprule
Statement & From \\
\midrule
Total baryon number and total charge vanish across both sectors & \Sig{0} \\
The coupling $\QCD$ is non-zero for a topological reason (precise form and proof) & \Sig{4}; used by \Max{2}, \Max{3} \\
The two homology generators of the torus as the origin of exactly two transparent fields & \Sig{2}, \Sig{T} \\
Punctures directed inwards & \Sig{P} \\
The wall picture and the four sectors re-expressed in terms of $\tau$ & \Sig{3} \\
\bottomrule
\end{tabular}
\end{center}
\end{table}

Table~\ref{tab:unowned} lists five statements of the earlier axioms that no paper of the series is yet assigned to. They are targets, not premises.

\section{The fifth axiom}\label{sec:axiom}

\subsection{What the four axioms do not reach}\label{sec:reach}

\Sig{5} and \Max{2} say the same thing: every apparent singularity in a lower-dimensional description is a projection artefact of a regular object on the full geometry, and the completed object is unique. Three assertions are bundled in it. The first is that what is observed is the image of something on $\Mtwelve$ under a passage to fewer dimensions. The second is that this passage loses information. The third is that for everything observed, including what is singular, a regular object above it exists, and that it is unique.

The four axioms and M1 and M2 supply the first of these as a scheme, since M2 obtains four-dimensional statements by integrating over $\Keight$ (M2 Definition~2.5). The second is a remark with a proof, given below. The third is not a consequence of the four axioms: they determine the geometry and the field equation, and they say nothing about which four-dimensional configurations arise. Its existence half is what the earlier work uses, and we declare it. Its uniqueness half we do not declare: it is not needed for the statement, the passage down is many-to-one, and in the series the object is fixed by the four axioms and minimality (M1 Theorem~8.1) and not by its shadow.

\subsection{The axiom}\label{sec:fifth}

\begin{axiom}[shadow]\label{ax:A5}
We can see the shadow, but the shadow is missing information about what it is a shadow of.
\end{axiom}

To use the sentence we must say what a shadow is, as M1 did for each of the four.

\begin{definition}[configurations and the reduction]\label{def:reduction}
A \emph{configuration} is a set of the fields of M2 Definition~2.5 on $\Mtwelve$ --- the metric with its warp, the dilaton $\QCD$, the $\Spinc$ spinor and the connection of the determinant line --- satisfying the field equation~\eqref{eq:master}. It is \emph{regular} if every field is defined at every point of $\Mtwelve$ and is smooth away from the corners and walls of the pillowcase, where M1 admits cone points and M2 admits a jump in the slope of the warp. The \emph{reduction} $\rho$ of a configuration is the four-dimensional data on $S^{1}\times\Sigma$ obtained by integrating over $\Keight$ in the two stages of M2 Definition~2.5. An \emph{observed configuration} is four-dimensional data of the same kind, defined on $S^{1}\times\Sigma$ away from a closed set of measure zero, on which it may fail to be defined or to be bounded.
\end{definition}

\begin{definition}[shadow, formal]\label{def:A5}
The geometry satisfies A5 if every observed configuration is the reduction of a regular configuration.
\end{definition}

The definition is an existence statement. It does not say that the regular configuration is unique, and it does not name a mechanism by which it is found. It is the minimal option among those the earlier wording leaves open, and is recorded as such: the earlier axiom declared existence, uniqueness and two named mechanisms; we declare existence only. \rigD\ (minimality)

\begin{remark}[the reduction loses information]\label{rem:loss}
The reduction is not one-to-one on fields. Integration over $\Keight$ sends to zero every function on $\Keight$ of zero mean, so two sets of fields that differ by such a function in one component have the same four-dimensional data. \rigR\ This is a statement about the map on fields. Whether two configurations, each satisfying~\eqref{eq:master}, can have the same reduction is a question about the solutions of the field equation, which this paper does not answer; the axiom does not depend on the answer.
\end{remark}

\subsection{The smallest case by hand}\label{sec:smallest}

The pillowcase is the smallest object of the series that has a singular point and a regular object above it, and every part of the axiom can be seen in it.

\begin{proposition}[the corner of the pillowcase]\label{prop:smallest}
Let $P=\Ttwo/\Ztwo$ be the pillowcase and $\pi:\Ttwo\to P$ the fold.
\begin{enumerate}[label=(\roman*),nosep]
\item $P$ has four cone points of angle $\pi$, at which all of its curvature is concentrated; each is the image under $\pi$ of a point of the torus at which the torus is flat and regular.
\item The fold is two-to-one away from the four corners and identifies $z$ with $-z$: it is not one-to-one.
\item The torus is not the only regular surface lying over $P$. For every $n\ge2$ the torus $\R^{2}/n\Z^{2}$ maps onto $\Ttwo=\R^{2}/\Z^{2}$, and hence onto $P$, by a map that commutes with the fold.
\end{enumerate}
\end{proposition}

\begin{proof}
(i) is M1 Theorem~5.7 and the paragraph that follows it there: near a fixed point the fold is the rotation by $\pi$, a disc of angle $2\pi$ is folded to a cone of angle $\pi$, and the torus is flat. (ii) The fold sends $z$ and $-z$ to the same point, and $z\neq-z$ except at the four half-lattice points. (iii) The identity of $\R^{2}$ descends to a map $\R^{2}/n\Z^{2}\to\R^{2}/\Z^{2}$, since $n\Z^{2}\subset\Z^{2}$, and it commutes with $z\mapsto-z$; composing with $\pi$ gives a map onto $P$. The finite content is certified on the grids $(\Z/4)^{2}$ and $(\Z/8)^{2}$: the fold is an involution with exactly four fixed points, the other twelve points fall into six pairs, two distinct points have the same image, every point of the smaller grid has four preimages in the larger, and the reduction commutes with the fold (\texttt{L1\_fold\_involution}, \texttt{L1\_four\_fixed}, \texttt{L1\_two\_to\_one}, \texttt{L1\_not\_injective}, \texttt{L1\_second\_cover}, \texttt{L1\_second\_cover\_fold}). \rigR
\end{proof}

In the words of the axiom: the corner is singular in the pillowcase and regular in the torus (existence); the pillowcase does not record which of $z$ and $-z$ a point came from (loss of information); and the pillowcase alone does not say which torus was folded (no uniqueness). The proposition is an illustration in dimension two and nothing later rests on it.

\subsection{What justifies declaring it}\label{sec:justify}

The earlier work offered a justification, \Max{4}: that the statement cannot be proved within any single existing formal system, three approaches having each reached an obstruction. We keep the observation and limit what is claimed for it. That three approaches did not close is evidence that the statement is hard; it is not a proof that the statement is undecidable, and nothing in the earlier text exhibits it as independent of a named formal system. We therefore argue, and do not prove, that the statement is to be taken as an axiom, on the same ground as the four: it is an observation in plain words whose content the others do not supply. \rigM

The axiom has a failure test, as every minimal choice of the series has. It fails if an observed configuration is shown to have no regular configuration above it. It is too weak if a later paper needs the uniqueness of the regular configuration, or a named mechanism of completion, as a premise and not as a result; in that case the choice made here is to be flagged and enlarged in the open. The first paper that tests it on a case beyond the pillowcase is M5.

\section{What is not an axiom}\label{sec:notaxioms}

Three of the twenty are of class N.

\Max{3} states that each proof follows one template: identify the shadow, construct the completion, prove a lock, derive the consequence. This is a method. It is not a statement about the universe or about the mathematics, and a method is adopted, not assumed. Two of its four steps rest on statements that are not in hand: the completion, which is the fifth axiom, and the lock, which is the rider of \Sig{4} listed in Table~\ref{tab:unowned}.

\Max{4} states that the Shadow Principle cannot be proved from within. It is a statement about provability and not an axiom of the theory. Its use is confined to Section~\ref{sec:justify}.

\Sig{7} states that complexity is a projection artefact of simplicity. It has no formal statement, no result of the earlier work cites it as a premise, and neither the definition of minimality nor the theorem of the least dimension in M1 asserts it. It is an interpretation, close in spirit to the fifth axiom, and is recorded as such.

Three clauses of other axioms are of the same kind and are not counted separately: the clause of \Sig{2} on the undecidability of other geometries; the clause of \Sig{T} that structured diversity requires two transparent fields, which is an argument; and the choice of unit in the earlier treatment of the warp, which M2 carries as a statement of units. No axiom is a convention as a whole; the conventions of the series are the assignment of two corner labels (M1 Appendix~B), which sheet is called $\psi$, and the overall sign of $\psi$ (M2).

\section{The minimal set, and the list of targets}\label{sec:result}

\begin{theorem}[the minimal set]\label{thm:minimal}
The series rests on the four axioms of M1, the fifth axiom of this paper, the principle of minimality, and one measured scale; of the twenty axioms of the earlier statement none other is independent.
\end{theorem}

\begin{proof}
By Theorem~\ref{thm:classification}, eight of the twenty are theorems of M1 and M2, which rest on the four axioms, the five readings and minimality; seven are targets, which are not premises; three are not axioms; and the remaining two are the one statement declared as A5 (\texttt{L2\_one\_axiom}, \texttt{L2\_minimal\_set}). The scale is the declared input of M1 Theorem~8.1 (\texttt{L2\_inputs}). \rigD
\end{proof}

The theorem is to be read with its limits. It says that no further axiom is needed to state what M1 and M2 have proved and what the later papers are to prove. It does not say that the targets of Table~\ref{tab:targets} will be reached from five axioms; if one of them proves not to be derivable it returns to the list of axioms, with the failure recorded. The five readings of M1 and M2 are definitions and not axioms, and a reader who takes one of them differently has a different theorem. And the count of ``one measured scale'' is of dimensionful inputs; the claim that there is no dimensionless one rests on the stationarity conditions of M1 Proposition~8.2, which M7 owes.

Tables~\ref{tab:targets} and~\ref{tab:unowned} are the list of what remains. They are the content of the papers that follow: the spectrum and the partition function, the test of the fifth axiom on the pillowcase, the aspect ratio, the masses and mixings, the vacuum energy, the black-hole interior and the couplings.

\section*{Acknowledgements}
The author used Claude (Anthropic) and Gemini (Google DeepMind) as mathematical development and analytical tools during the preparation of this work. All intellectual content, theoretical framework, and scientific claims are the sole responsibility of the author.

\appendix

%%=====================================================================
%% M3 Appendix A -- the twenty axioms, transcribed as data (SOW node I.1).
%% GENERATED from sources/Paper_XX_New_Axioms_for_Old_Problems_v10.tex,
%% ll.335-555 (Sigma) and ll.639-770 (the eight axiom environments only).
%% Each block carries its source line range. Departures from the source,
%% all by PI ruling of 10 Nov 2026 (WP0 assumptions 4-6):
%%   (a) every reference to a paper of the earlier corpus is omitted and
%%       marked [\ldots]; external citations are kept;
%%   (b) the three places where omega is named in a superseded form carry
%%       the canonical name in square brackets;
%%   (c) the source's local macros are expanded; the rendered text is the same;
%%   (d) layout only: headings are unnumbered; the Sigma-3 table is not a float;
%%       the eight mathematical axioms are set under headings, not in an
%%       axiom environment; the labels of the source are dropped;
%%   (e) not transcribed: the Example paragraphs after each mathematical
%%       axiom, and the Fixed-Point Vanishing lemma block (ll.659-718).
%%   (f) five paragraphs headed "Note added in M3" are not source text: the
%%       PI's commentary of 10 Nov 2026 (Sigma-1, Sigma-2, Sigma-2f), the
%%       correction of the Lyapunov exponent in M-7 (P17), and the correction
%%       of d_odd in M-6 (number, not sum, of odd divisors).
%%   (g) layout only: the quoted text is set with extra line-breaking
%%       tolerance (the group opened after the source note and closed at the
%%       end of the file), so that no line overruns the text block.
%% External citation keys needed in M3's bibliography are listed at the end.
%%=====================================================================

\section{The twenty axioms, as stated}\label{app:axioms}

% PROVISIONAL source note (wording for the PI; no M3 prose is settled yet).
\noindent The twenty axioms audited in this paper are reproduced here as data.
The wording is that of the programme's earlier statement of them; that
document is superseded by the present series and is not cited. Line numbers
refer to its tenth version. An omission is marked $[\ldots]$; words in square
brackets replace a superseded name of~$\omega$. A paragraph headed \emph{Note
added in M3} is commentary of the present paper and is not part of the earlier
wording. The twelve creation axioms
come first ($\Sigma$-3a is a part of $\Sigma$-3), then the eight mathematical
axioms.

\begingroup\tolerance=3000\emergencystretch=4em


\subsection*{The creation axioms}

% ---- source ll.335-345 ----
\subsection*{$\Sigma$-0: Universal Symmetry}
\noindent{\small Source lines 335--345.}\par\smallskip
\textbf{The universe is exactly symmetric across the $\Ztwo$ involution.}

Every apparent asymmetry---matter/antimatter, L/D chirality, baryon excess, visible/dark---is a confinement artefact of observation from one wall of the pillowcase.  The totality is perfectly balanced.  Total baryon number across both walls: zero.  Total chirality across both walls: zero.  Total charge across both walls: zero.

\emph{Plain English: The universe is perfectly balanced.  Everything that looks lopsided---more matter than antimatter, left-handed amino acids, visible matter versus dark matter---looks that way only because we are seeing one side.  The other side has the mirror image.  Together, perfect symmetry.}

\smallskip

\noindent\emph{Mathematical lineage.}\ Universal Symmetry is the axiomatic statement of a century of physics-led reasoning about what symmetry demands.  Dirac (1928, \emph{The Quantum Theory of the Electron})~\cite{Dirac1928} derived the relativistic wave equation for the electron and, confronting its negative-energy solutions, reasoned that symmetry required a mirror particle --- the positron, confirmed by experiment four years later.  L\"{u}ders (1954)~\cite{Luders1954} and Pauli (1955)~\cite{Pauli1955} independently proved the CPT theorem: every Lorentz-invariant local quantum field theory is invariant under the combined operation of charge conjugation, parity, and time reversal --- the deepest symmetry statement in physics, asserting that apparent asymmetries under any one operation must be balanced by the others.  Philip~W.\ Anderson (1972, ``More is Different'')~\cite{Anderson1972} established broken symmetry as observational fact: the ground state of a many-body system need not respect the symmetry of its underlying laws, and what the observer measures is the state, not the law.  $\Sigma$-0 extends this lineage by identifying the missing wall: apparent asymmetries --- matter/antimatter, chirality, baryon excess, visible/dark --- are not consequences of broken laws but confinement artefacts of observation from one wall of the $\Ztwo$ pillowcase. [\ldots]  The axiom therefore names the principle; it is an identification, not an invention.

% ---- source ll.348-360 ----
\subsection*{$\Sigma$-1: The Manifold}
\noindent{\small Source lines 348--360.}\par\smallskip
\textbf{We are four-dimensional entities living in a twelve-dimensional universe.  The additional dimensions are vast, not small.}

What we observe in four dimensions is a projection, and projections lose information.  Every phenomenon that appears singular, paradoxical, or fine-tuned in the projection is smooth and natural in the full geometry.  This is not the string-theory compactification paradigm of dimensions curled into Planck-scale tubes.  The dimensions are vast.

\noindent\emph{Note added in M3.}\ A picture for the reader: the universe is like a set of nested Russian dolls. We are in the fourth doll, and we are contained within eight other, larger dolls, each larger than the one before.

Formally: $\Mtwelve = (S^1_{\mathrm{Sol}} \times S^3) \times K_8$, where $K_8 = [(\mathbb{CP}^2 \times \Stwo) \times_w (\Ttwo/\Ztwo)]^{\Spinc}$, with Euler characteristic $\chi(K_8) = 12$.

\emph{Plain English: We are four-dimensional beings in a twelve-dimensional universe.  Those extra dimensions are not tiny curled-up strings---they are vast.  We are seeing a shadow of the full reality, and the shadow loses information.}

\smallskip

\noindent\emph{Mathematical lineage.}\ The Manifold is the axiomatic statement of a 200-year tradition that treats the dimensionality and intrinsic geometry of space as empirical rather than a priori.  Gauss (1827, \emph{Disquisitiones generales circa superficies curvas})~\cite{Gauss1827} established the Theorema Egregium: curvature is intrinsic --- geometry need not depend on embedding, and can be known from measurements within the manifold alone.  Riemann (1854, \emph{\"{U}ber die Hypothesen welche der Geometrie zu Grunde liegen}, his G\"{o}ttingen Habilitationsvortrag)~\cite{Riemann1868} generalised Gauss's intrinsic geometry from two-dimensional surfaces to $n$-dimensional manifolds, reframing dimension itself from a priori fact to empirical question.  Kaluza (1921, ``Zum Unit\"{a}tsproblem der Physik'')~\cite{Kaluza1921} executed the first physical realisation: Einstein's equations in five dimensions automatically produce gravitation and electromagnetism in four, unifying the two forces through a single geometric structure.  $\Sigma$-1 extends the lineage: the twelve dimensions of $\Mtwelve$ are vast and topologically isolated, a direct realisation of Riemann's empirical dimensionality. [\ldots]  The axiom therefore names the principle; it is an identification, not an invention.

% ---- source ll.363-427 ----
\subsection*{$\Sigma$-2: The Fundamental Shapes}
\noindent{\small Source lines 363--427.}\par\smallskip
\textbf{The twelve-dimensional manifold consists of specific geometric components.  Each is fundamental and irreducible.  Given these shapes, everything follows.}

There are no free parameters within this geometry.  The Master Equation, the Standard Model, the fine structure constant, the cosmological constant, dark matter, consciousness, and the spectral structure of the primes are all consequences of these specific shapes and their interactions.  Each component is independently necessary: remove any one and the framework collapses.

Whether alternative self-consistent geometries could produce other universes is undecidable---not merely from within the MEF as a formal system, but from within this universe.  The shapes are not a description of reality that we could in principle examine from outside.  They \emph{are} reality.  We are inside them.  This is the \textbf{External Horizon} of the G\"odel Boundary.

\noindent\emph{Note added in M3.}\ By M1 (Theorem~8.1, the minimality theorem) this is the minimal topology that satisfies the four axioms. There may well be higher dimensions of which this universe is a subset. That is unknowable from where we are.

\emph{Plain English: The universe is built from exactly these shapes.  Each is specific and each is necessary.  Everything we observe---every force, every particle, every constant---drops out of these shapes and nothing else.  Could different shapes have made a different universe?  We cannot know.  Not because we haven't worked it out yet, but because we are inside these shapes---they are our reality, not a model we look at from outside.}

\subsubsection*{$\Sigma$-2a: $\mathbb{CP}^2$---The Colour Engine}

Complex projective 2-space.  Isometry group $= \mathrm{SU}(3)$: this is the strong force.  Dolbeault cohomology $H^{0,2}(\mathbb{CP}^2, \mathcal{O}(-4))$ gives exactly three fermion generations.  $w_2(\mathbb{CP}^2) \neq 0$ forces $\Spinc$ ($\Sigma$-2f).  $c_1(\mathcal{O}(3))^2 = 9$ feeds into the topological weight $W = 36$ that determines the fine structure constant.

\emph{Remove $\mathbb{CP}^2$ and you lose: the strong force, three generations, all quark physics, the fine structure constant.}

\subsubsection*{$\Sigma$-2b: $\Stwo$---The Electroweak Sphere}

The unique simply-connected 2-manifold with $\mathrm{SU}(2)$ isometry.  $\chi(\Stwo) = 2$.  Electroweak breaking gives $\sin^2\thetaW = 0.224$.  Geometric tension versus $\Spinc$ monopole flux produces the Higgs mechanism.

\emph{Remove $\Stwo$ and you lose: the weak force, electroweak symmetry breaking, the Weinberg angle, the Higgs mechanism.}

\subsubsection*{$\Sigma$-2c: $\Ttwo$---The Torus}

Two independent homology generators are the topological origin of the two transparent fields (gravity and $\omega$).  Doubly-periodic structure enables modular forms.  Aspect ratio $\tau_2 = L_\omega/L_\psi = 1.86$ is fixed by $\Spinc$ flux stabilisation, not a free parameter.

\emph{Remove $\Ttwo$ and you lose: modular forms, the partition function, the entire spectral structure, the $\omega$-axis.}

\subsubsection*{$\Sigma$-2d: $\Ztwo$---The Involution}

The order-2 group acting on $\Ttwo$ by $z \mapsto -z$, creating the orbifold $\Ttwo/\Ztwo$ (the pillowcase).  Four fixed points define the four sectors of reality ($\Sigma$-3).  Q-wall/AQ-wall separation is forced.  Generates mock modularity.  Creates all apparent asymmetry.

\emph{Remove $\Ztwo$ and you lose: the shadow sector, dark matter, chirality, mock modularity, the four-sector pillowcase structure.}

\subsubsection*{$\Sigma$-2e: $\times_w$---The Warped Product}

The fibration connecting $\mathbb{CP}^2 \times \Stwo$ to $\Ttwo/\Ztwo$ is warped, not direct.  Creates the Planck--electroweak hierarchy.  Enables $\Lambda_{\mathrm{eff}} \approx 10^{-122}\,M_{\mathrm{Pl}}^4$.

\emph{Remove $\times_w$ and you lose: the hierarchy between fundamental forces, the cosmological constant explanation, the entire scale structure of physics.}

\subsubsection*{$\Sigma$-2f: $\Spinc$---The Twist}

The $\Spinc$ structure on $K_8$, forced by $\mathbb{CP}^2$.  Enables fermions.  Breaks Liu rigidity.  Determinant line bundle $L = \mathcal{O}(3)$ gives $\mathrm{U}(1)_Y$.  Simultaneously creates the disease (mock modularity) and provides the cure (Zwegers completion).

\emph{Remove $\Spinc$ and you lose: all fermions, all dynamics, hypercharge.  The universe cannot exist.}

\emph{Plain English: The pillowcase has a spin that creates rigidity.  That rigidity is broken through asymmetry---the $\Spinc$ twist.  Without it, nothing would happen.  The universe would be frozen.  The twist is what makes everything move, evolve, live.}

\noindent\emph{Note added in M3.}\ The same condition says what the next time slice is to be, so that each time slice is different. In the series this is the content of A3 (each instance is unique) and of A4 (the universe never stands still), with the twist as their carrier. Time can move forward; without this condition everything would be frozen.

\subsubsection*{$\Sigma$-2g: $S^1_{\mathrm{Sol}}$---The Temporal Circle}

The circle carrying the Sol monodromy $M = \bigl(\begin{smallmatrix} 2 & 1 \\ 1 & 1 \end{smallmatrix}\bigr) \in \mathrm{SL}(2,\mathbb{Z})$.  Eigenvalues $\phi^2$ and $\phi^{-2}$.  Anosov property generates mixing, chaos, the arrow of time.  $D(\tilde{M}) = 5 = \mathrm{disc}(\Q(\sqrt{5}))$.

\emph{Remove $S^1_{\mathrm{Sol}}$ and you lose: the arrow of time, mixing, chaos, Anosov dynamics, Fibonacci growth, the cascade pump in turbulence, the conservation law $F(\tfrac{1}{2}) = 1$.}

\emph{Plain English: The Sol monodromy is the clock of the universe.  It ticks at the golden ratio.  Every new growth---in time, in nature, in a spiral---is displaced by $\phi$.  That is why we see Fibonacci everywhere, we live in a fractal universe.}

\subsubsection*{$\Sigma$-2h: $S^3$---The Spatial Sphere}

The unique simply-connected closed 3-manifold with constant positive curvature (Perelman's theorem, resolving the Poincar\'{e} Conjecture).  Closed, finite spatial topology.  Discrete spectrum.  CMB low-$\ell$ suppression.

\emph{Remove $S^3$ and you lose: spatial closure, discrete spectrum, CMB low-$\ell$ suppression, closed-universe predictions.}

\subsubsection*{The Euler Characteristic}

$\chi(K_8) = \chi(\mathbb{CP}^2) \times \chi(\Stwo) \times \chi(\Ttwo/\Ztwo) = 3 \times 2 \times 2 = \mathbf{12}$.  This single topological integer propagates through the entire framework: fermion generations $= \chi/4 = 3$; coupling $\zeta = \pi\chi = 12\pi$; hierarchy $e^{12\pi}$; cosmological constant exponent $e^{-96\pi}$ where $96 = 8 \times 12$; BAO spectral bridge $e^{-\zeta/\dim(K_8)} = e^{-3\pi/2}$.

% ---- source ll.430-449 ----
\subsection*{$\Sigma$-3: The Pillowcase}
\noindent{\small Source lines 430--449.}\par\smallskip
\textbf{The control plane ($\Ttwo/\Ztwo$) has two walls and a drive axis, creating four sectors.}

The Quantum Wall ($\psi$): the visible, measurable sector.  The Anti-Quantum Wall ($\bar{\psi}$): the shadow sector---structurally essential, not directly observable from the Q-wall.  The $\omega$-axis: [the Coherent Field], from low (ordered, classical) to high (interconnected, quantum).

\begin{center}
\small
\begin{tabular}{@{}l p{5cm} p{5cm}@{}}
\toprule
& \textbf{Low $\omega$ (ordered)} & \textbf{High $\omega$ (diverse)} \\
\midrule
\textbf{$\psi$ (Q-wall)} & Standard Model, visible matter, classical physics & Life, consciousness, creativity, negentropy \\[3pt]
\textbf{$\bar{\psi}$ (AQ-wall)} & Dark matter---gravitationally present, electromagnetically invisible & Dark consciousness---the shadow at maximum coherence \\
\bottomrule
\end{tabular}
\end{center}

\emph{Plain English: The pillowcase has four areas.  On the quantum wall: ordinary matter at low coherence, and living conscious beings at high coherence.  On the anti-quantum wall: dark matter at low coherence, and dark consciousness at high coherence.}

% ---- source ll.452-473 ----
\subsection*{$\Sigma$-3a: The Quaternionic Layer}
\noindent{\small Source lines 452--473.}\par\smallskip
\textbf{The fibre algebra of $K_8$ is non-abelian over an abelian base, and the non-commutativity resolves to the classical real limit at exactly one of the four corners.}

The pillowcase $\Ttwo/\Ztwo$ carries two structurally distinct algebraic layers.  The \emph{base} layer is the Klein-four group $V_{\mathrm{geom}} = (\Z/2)^2$ of half-period translations of $\Ttwo$, whose four elements index the four corners; together with the centre $Z = \langle -1 \rangle$ of the quaternion group, the base structure $V_{\mathrm{geom}} \times Z$ is abelian.  The \emph{fibre} layer is the quaternion algebra $\mathbb{H}$ acting on the pillowcase interior, with unit-quaternion group $Q_8 = \{\pm 1, \pm i, \pm j, \pm k\}$ --- the 8-point non-abelian multiplicative group.  Non-abelian fibre, abelian base.

At each corner of the pillowcase the fibre algebra degenerates to a sub-algebra of $\mathbb{H}$.  The four limits are
\[
  C_0 \longrightarrow \R, \quad
  C_i \longrightarrow \C_i, \quad
  C_j \longrightarrow \C_j, \quad
  C_k \longrightarrow \C_k,
\]
with dimensions $1, 2, 2, 2$ respectively.  Only at $C_0$ does the fibre reach the fully-commutative classical limit $\R$; the other three corners retain a non-trivial imaginary unit and remain complex.  The four corners are therefore not symmetric under the $\mathrm{Aut}(\mathbb{H})$ action: $C_0$ is uniquely distinguished as the corner where the non-commutativity of the interior commutativises completely.

This asymmetry is the structural source of the three Millennium closures on $\Mtwelve$.  The Riemann critical line $\sigma = \tfrac{1}{2}$ is the fixed locus of the transpose anti-involution on the hyperfinite completion of the fibre algebra, and that fixed locus intersects $\mathbb{H}$ at exactly $\R$ --- i.e., at $C_0$ [\ldots].  The Yang--Mills mass gap is the minimum geodesic distance on $K_8$ between distinct corners under the $\mathrm{Sp}(1)$-action on the fibre; the alternation $K_2(n, -1; 2) = (-1)^n$ is the orbit signature of $V_{\mathrm{Aut}} \times Z$ on $Q_8$ over $n$ topological wraps of the base [\ldots].  The Navier--Stokes velocity field embeds into $\mathbb{H}$ as a pure-imaginary quaternion; the Zwegers completion $\hat f = f + g^*$ is the algebraic closure $\mathrm{Im}(\mathbb{H}) \hookrightarrow \mathbb{H}$ that re-admits the $\R$-component generated by Hamilton multiplication [\ldots].  Each of the three closures uses a different structural feature of this single non-abelian-over-abelian algebra; the uniqueness of $C_0$ is what selects the Riemann critical line, the Yang--Mills vacuum corner, and the Navier--Stokes scalar-absorption channel.

\emph{Plain English: The pillowcase has a hidden algebraic structure that is deeper than the four sectors of $\Sigma$-3.  The interior carries Hamilton's quaternions, where rotations do not commute --- order matters.  The base, which labels where you are, is commutative --- order does not matter.  At three of the four corners, the quaternion structure simplifies to ordinary complex numbers, keeping a trace of its quaternionic parent.  At exactly one corner --- the origin --- it collapses all the way to the real numbers, fully commutative.  That one unique corner is what makes the three Millennium problems solvable: it is where the Riemann critical line lives, where the Yang--Mills vacuum sits, and where the Navier--Stokes real-component absorption happens.  The asymmetry of the four corners is not decorative.  It is the load-bearing structural feature of the manifold.}

\smallskip

\noindent\emph{Mathematical lineage.}\ The Pillowcase is the axiomatic statement of a two-century tradition in non-commutative algebra and its classification.  Hamilton (1843, Brougham Bridge)~\cite{Hamilton1844} constructed the quaternions $\mathbb{H}$ as the first non-commutative number system supporting geometric rotation --- the only honest step beyond $\C$ for half a century.  Jordan (1870, \emph{Trait\'{e} des substitutions et des \'{e}quations alg\'{e}briques})~\cite{Jordan1870} established the systematic theory of finite groups, whose classification methods and structural insights underlie all subsequent work on discrete symmetries --- including the four-fixed-point structure of the pillowcase base.  Frobenius (1877, \emph{\"{U}ber lineare Substitutionen und bilineare Formen})~\cite{Frobenius1877} closed the division-algebra classification: $\R$, $\C$, $\mathbb{H}$ are the only finite-dimensional associative division algebras over $\R$; the octonions $\mathbb{O}$ extend the hierarchy non-associatively and do not feature in the pillowcase's algebraic spine.  $\Sigma$-3 realises the associative sub-hierarchy $\R \subset \C \subset \mathbb{H}$ geometrically over the four corners of $\Ttwo/\Ztwo$: $\mathbb{H}$ in the interior of the fibre (non-commutative), $\C$ at three corners (complex structure preserved), $\R$ at the unique corner $C_0$ (fully commutative).  The pillowcase is the Frobenius hierarchy carrying geometry. [\ldots]  The axiom therefore names the principle; it is an identification, not an invention.

% ---- source ll.476-482 ----
\subsection*{$\Sigma$-4: The Three Fields}
\noindent{\small Source lines 476--482.}\par\smallskip
\textbf{Three fundamental geometric degrees of freedom.  Their interaction---not any one alone---produces the coupling.}

$\omega$: [the Coherent Field].  $\psi$: quantum spinor condensate (Q-wall).  $\bar{\psi}$: anti-quantum spinor condensate (AQ-wall).  The derived coupling: $\QCD = \omega(\psi - \bar{\psi})^2$.  $\QCD \neq 0$ topologically: merging the walls would destroy the $\Ztwo$ orbifold.

\emph{Plain English: Three fields.  $\omega$ is [the Coherent Field].  $\psi$ and $\bar{\psi}$ are chiral mirror partners living on opposite walls.  The coupling constant $\QCD$ measures how the engine operates, as well as the magnitude of its effect.}

% ---- source ll.485-495 ----
\subsection*{$\Sigma$-5: The Shadow Principle}
\noindent{\small Source lines 485--495.}\par\smallskip
\textbf{Every apparent singularity in lower dimensions is a projection artefact of a smooth, complete object in higher dimensions.  There are no singularities in the universe.}

What looks singular in $n$ dimensions is smooth in $n+1$.  A blow-up in 3D fluid mechanics is the shadow of a smooth 4D object.  An off-line zero of $\zeta(s)$ in the 2D complex plane is the shadow of a smooth 3D spectral object.

\emph{Plain English: There are no singularities in the universe.  Every time we see something that looks like a singularity---a black hole, a blow-up, an impossible zero---we are seeing the shadow.  Go up one dimension and it's smooth.}

\smallskip

\noindent\emph{Mathematical lineage.}\ The Shadow Principle is the axiomatic statement of a two-hundred-year tradition.  Jacobi (1829, \emph{Fundamenta Nova})~\cite{Jacobi1829} showed that abelian integrals, multivalued on a Riemann surface, become single-valued on the Jacobian torus --- the completion lifts pathological shadow behaviour to smooth structure one dimension up.  Ramanujan (1920, deathbed letter to Hardy)~\cite{Ramanujan1920} announced seventeen mock theta functions that behaved like modular forms but were not, with no explanation of what ``mock'' meant.  Zwegers (2002, Utrecht thesis)~\cite{Zwegers2002} resolved the eighty-two-year mystery by constructing the non-holomorphic completion $\hat{f} = f + g^{*}$, where $g^{*}$ is the Eichler integral of the shadow $g$ --- a direct generalisation of Jacobi's theta machinery to half-integer weight and congruence subgroups.  $\Sigma$-5 promotes this mechanism to a universal axiom: Jacobi's periods, Ramanujan's mock theta functions, and Zwegers' completions are not three separate phenomena but three surfaces of one structural principle [\ldots].  The axiom therefore names the principle; it is an identification, not an invention.

% ---- source ll.498-508 ----
\subsection*{$\Sigma$-6: Shape Determinism}
\noindent{\small Source lines 498--508.}\par\smallskip
\textbf{Everything is fixed by the shapes.  There are no free parameters.}

Every physical quantity---the fine structure constant, the baryon asymmetry, the mass ratios, the cosmological constant, the dark-to-visible matter ratio---is a topological invariant of $\Mtwelve$ or a geometric consequence of $\chi(K_8) = 12$.

\emph{Plain English: Everything is what it is because of the shape of the 12-dimensional manifold.  No knobs, no tuning, no choices and no mysticism.}

\smallskip

\noindent\emph{Mathematical lineage.}\ Shape Determinism is the axiomatic statement of a 125-year programme to derive physics from mathematics with no free parameters.  Hilbert (1900, Paris ICM, Sixth Problem)~\cite{Hilbert1900} posed the foundational challenge: axiomatise the physical sciences on the model of geometry.  Einstein (1915, \emph{Die Feldgleichungen der Gravitation})~\cite{Einstein1915} executed the first complete instance --- general relativity derives gravitation from the curvature of spacetime, with no free parameters beyond Newton's constant required for the Newtonian limit.  Weyl (1918, \emph{Gravitation und Elektrizit\"{a}t})~\cite{Weyl1918} attempted the first extension: a geometric derivation of electromagnetism through a gauge principle.  The 1918 theory failed on Einstein's clock-rate objection, but the gauge principle itself became the foundation of all subsequent fundamental physics.  $\Sigma$-6 claims the extension can be completed: every physical quantity --- the fine structure constant, mass ratios, the cosmological constant, the dark-to-visible ratio --- is a topological invariant of $\Mtwelve$ or a geometric consequence of $\chi(K_8) = 12$. [\ldots]  The axiom therefore names the principle; it is an identification, not an invention.

% ---- source ll.511-517 ----
\subsection*{$\Sigma$-7: Dimensional Simplicity}
\noindent{\small Source lines 511--517.}\par\smallskip
\textbf{Complexity is a projection artefact of simplicity.}

Things look complex in lower dimensions because we are seeing the second, third, fourth, or fifth derivative of what is a much simpler shape in higher dimensions.

\emph{Plain English: Incredibly complex things look complex because we are looking at the derivative of what is a much simpler shape in the higher dimension.  Go up, and complexity dissolves into simplicity.}

% ---- source ll.520-526 ----
\subsection*{$\Sigma$-8: The $\omega$-Drive}
\noindent{\small Source lines 520--526.}\par\smallskip
\textbf{Low $\omega$: discrete, disconnected, classical.  High $\omega$: interconnected, quantum, diverse.}

At low $\omega$, systems look and behave as if they are independent.  At high $\omega$, interconnectedness intensifies.  What appears as chaos at high $\omega$ is structured diversity too rich for low-$\omega$ observation to parse.

\emph{Plain English: Low coherence: things are separate, classical, simple.  High coherence: everything is connected, quantum effects dominate, and what looks like chaos is actually structure too rich for our instruments to parse.}

% ---- source ll.529-535 ----
\subsection*{$\Sigma$-P: Dimensional Permeability}
\noindent{\small Source lines 529--535.}\par\smallskip
\textbf{The boundaries between dimensional shells admit three types of interaction.  These are exhaustive.}

\textbf{(i)~Transparent.}  The field permeates all twelve dimensions.  Gravity and $\omega$ are the only transparent fields ($\Sigma$-T).  \textbf{(ii)~Puncture.}  Topology is broken, creating a conduit.  Black holes are outward conduits; quasars are inward conduits.  \textbf{(iii)~Impermeable.}  Gauge forces cannot cross between shells.  Photons are imprisoned on their shell.  This converts perfect symmetry ($\Sigma$-0) into apparent asymmetry.

\emph{Plain English: Some things pass through all dimensions freely---gravity and $\omega$.  Some things punch through, creating holes.  And gauge forces are trapped.  That last one is why we think symmetry is broken when it isn't.}

% ---- source ll.538-544 ----
\subsection*{$\Sigma$-S: Shadow Influence}
\noindent{\small Source lines 538--544.}\par\smallskip
\textbf{The AQ-wall is not in a different dimension.  Shadow influence is structurally distinct from dimensional permeability.}

The AQ-wall is the chiral mirror of the Q-wall within the same geometry, separated by the $\Ztwo$ involution.  Its influence is felt through three channels: gravitationally (dark matter), through $\omega$ (dark consciousness), and through the mock modular completion (high-$\omega$ phenomena that appear inexplicable from the Q-wall alone: QCD vacuum condensate, superfluid behaviour, quantum tunnelling in enzymes, photosynthesis efficiency, prime conspiracies).

\emph{Plain English: The anti-quantum wall isn't in a different dimension---it's the mirror side of our geometry.  Its influence reaches us through gravity, through $\omega$, and through the mock modular mathematics.}

% ---- source ll.547-555 ----
\subsection*{$\Sigma$-T: Dual Transparency}
\noindent{\small Source lines 547--555.}\par\smallskip
\textbf{Structured diversity across boundaries requires at least two independent transparent fields.  The geometry provides exactly two.}

With one transparent field (gravity alone), every point beyond every boundary looks identical---just mass, no structure.  With two (gravity $+$ $\omega$), structure becomes resolvable.  Gravity reports \emph{how much}; $\omega$ reports \emph{how coherent}.  One gives detection.  Two give discrimination.  Two give the universe the capacity for structured diversity across every boundary.

The two homology generators of $\Ttwo$ ($\Sigma$-2c) are the topological origin of this requirement.  Two is the minimum for a universe of structured diversity.  Two is what the geometry provides.

\emph{Plain English: A universe with only one field crossing boundaries would be flat, dull, featureless.  You need two: a line and a branch.  Gravity is the line---it tells you something is there.  $\omega$ is the branch---it tells you what kind of thing it is.  One channel gives detection.  Two give navigation.}


\subsection*{The mathematical axioms}

% ---- source ll.639-641 ----
\subsubsection*{$\mathbb{M}$-1: Physical Observables as Internal Spectral Data}
\noindent{\small Source lines 639--641.}\par\smallskip
Physical observables---including the spectral zeta function $\zeta(s)^2(1-2^{-s})$---are internal spectral data of the $\Spinc$ CFT on $\Ttwo/\Ztwo$ with central charge $c = 12$.  They are output of the geometry, not external data imported into the theory.

% ---- source ll.647-649 ----
\subsubsection*{$\mathbb{M}$-2: The Shadow Completion}
\noindent{\small Source lines 647--649.}\par\smallskip
Every apparent singularity in the mathematical description---blow-up, off-line zero, massless state, divergent series---is a projection artefact.  The completed mathematical object on the full geometry of $\Mtwelve$ is smooth.  The completion is achieved by the Zwegers completion (in $\tau$-space) or the $\mathrm{GL}(3)$ Eisenstein completion (in $s$-space), and is unique given the topological data of $K_8$.

% ---- source ll.722-724 ----
\subsubsection*{$\mathbb{M}$-3: The Universal Proof Template}
\noindent{\small Source lines 722--724.}\par\smallskip
Each proof follows a single template: identify the shadow (incomplete projection) $\to$ construct the completion $\to$ prove the topological lock $\to$ derive the consequence.

% ---- source ll.730-732 ----
\subsubsection*{$\mathbb{M}$-4: The G\"{o}del Boundary (Internal Ceiling)}
\noindent{\small Source lines 730--732.}\par\smallskip
The Shadow Principle ($\Sigma$-5 / $\mathbb{M}$-2) cannot be proved from within any single existing formal system.  Three independent frameworks were tested---$\mathrm{GL}(3)$ spectral theory, the conformal bootstrap, transfer operators---and each reaches a structural wall.  This is not a failure of ingenuity.  It is G\"odel's incompleteness theorem in action: a true statement that cannot be proved from within the system.  The Shadow Principle is an axiom, justified by the totality of the MEF's consequences---just as Euclid's parallel postulate, Newton's $F = ma$, and Einstein's equivalence principle are axioms justified by their explanatory power.

% ---- source ll.740-742 ----
\subsubsection*{$\mathbb{M}$-5: The Partition Function}
\noindent{\small Source lines 740--742.}\par\smallskip
$Z_{K_8}^{\Spinc}(\tau,z)$ is a mock Jacobi form (weight~0, level~2, mock theta order~2).  Its completion $\hat{f} = f + \int g$ is a harmonic Maass form, smooth for $\operatorname{Im}(\tau) > 0$.

% ---- source ll.750-752 ----
\subsubsection*{$\mathbb{M}$-6: The Spectral Data}
\noindent{\small Source lines 750--752.}\par\smallskip
The eigenvalues of $D^2_{K_8}$ are $\lambda_n = n$ ($n \geq 1$).  The degeneracies are $m(n) = \dodd(n)$---the sum of odd divisors of $n$.  The spectral zeta function is $\zeta_{D^2}(s) = \zeta(s)^2(1-2^{-s})$.

\noindent\emph{Note added in M3.}\ Here $\dodd(n)$ is the number of odd divisors of $n$, not their sum: the Dirichlet series $\zeta(s)^{2}(1-2^{-s})$ stated in the axiom is that of the count, $\sum_{k\ \mathrm{odd}}q^{k}/(1-q^{k})=\sum_{m\geq1}\dodd(m)\,q^{m}$. For example $\dodd(6)=2$.

% ---- source ll.760-762 ----
\subsubsection*{$\mathbb{M}$-7: The Sol Monodromy}
\noindent{\small Source lines 760--762.}\par\smallskip
The monodromy of the $S^1$ factor is $M = \bigl(\begin{smallmatrix} 2 & 1 \\ 1 & 1 \end{smallmatrix}\bigr) \in \mathrm{SL}(2,\mathbb{Z})$, with eigenvalues $\phi^2$ and $\phi^{-2}$, where $\phi = (1+\sqrt{5})/2$ is the golden ratio.  The Lyapunov exponent is $\ln\phi$.  The fixed point of $\tau \mapsto (2\tau+1)/(\tau+1)$ is $\phi$.  The Anosov property generates mixing and the arrow of time.

\noindent\emph{Note added in M3.}\ The eigenvalues are as stated: $\phi^{2}$ and $\phi^{-2}$ (M1, Theorem~4.2; M2, certificate \texttt{L4\_eigen}). The Lyapunov exponent is therefore $\ln\phi^{2}=2\ln\phi$, not $\ln\phi$ as written above; M2 has $\kappa=\log\phi^{2}$.

% ---- source ll.768-770 ----
\subsubsection*{$\mathbb{M}$-8: The $\Spinc$ Twist}
\noindent{\small Source lines 768--770.}\par\smallskip
The $\Spinc$ structure on $K_8$---forced by $w_2(\mathbb{CP}^2) \neq 0$---breaks Liu rigidity, activating the $q$-dependent spectrum $\dodd(n)$.  Without it, the partition function would be Liu-rigid: $A_2 = 1 - y + y^2$, a topological invariant with no $q$-dependence.  The twist simultaneously creates the disease (mock modularity---apparent singularities, shadows, incompleteness) and provides the cure (the Zwegers completion---smoothness, the Shadow Principle).

\endgroup

%% External citation keys used above (entries to be carried into M3's bibliography):
%   Anderson1972
%   Dirac1928
%   Einstein1915
%   Frobenius1877
%   Gauss1827
%   Hamilton1844
%   Hilbert1900
%   Jacobi1829
%   Jordan1870
%   Kaluza1921
%   Luders1954
%   Pauli1955
%   Ramanujan1920
%   Riemann1868
%   Weyl1918
%   Zwegers2002


%%=====================================================================
%% M3 Appendix B -- cumulative symbol table: M1 and M2 reproduced, M3 added.
%% The M1 paragraph and the M2 rows are those of M2 Appendix B; the "where
%% defined" column gives the printed numbers of M2, since its labels are not
%% available here.
%%=====================================================================

\section{The labelling of the corners, and the symbols of the series}\label{app:labels}

\paragraph{From M1 (reproduced).}
The four corners are named by position, $P_{1}=(0,0)$, $P_{2}=(\tfrac12,0)$,
$P_{3}=(0,\tfrac12)$, $P_{4}=(\tfrac12,\tfrac12)$ in coordinates
$(\theta_{\psi},\theta_{\omega})$, and by algebra, $C_{0}=\R$, $C_{i}$,
$C_{j}$, $C_{k}$. Throughout the programme the dictionary is: $P_{1}=C_{0}$
(fixed by the drive; $\iota$ acts as conjugation); $P_{3}=C_{i}$ (the
$\omega$ half-period); $P_{2}=C_{j}$ (the $\psi$ half-period); $P_{4}=C_{k}$
(the corner at which the $\Spinc$ lift carries the sign $-1$). The assignment
of $j$ and $k$ between $P_{2}$ and $P_{3}$ is a convention and carries no
content (M1 Appendix~B).

\paragraph{Introduced in M2 (reproduced).}
\begin{center}\small
\begin{tabular}{@{}p{0.29\textwidth}p{0.49\textwidth}p{0.15\textwidth}@{}}
\toprule
Symbol & Meaning & Where defined \\
\midrule
$\psi,\ \bar\psi=\hat\tau\psi$ & the sheet field, a section of the Clifford module $M\cong\Hquat^{32}$, and its shadow & M2 Def.~2.2 \\
$M$, $\langle\cdot,\cdot\rangle_{M}$ & irreducible $\mathrm{Cl}(0,12)$-module; its inner product (trace) & M2 Def.~2.2, M2 \S3 \\
$\mathcal{A},\ \mathcal{C}$ & bundle of fibre algebras; corner-algebra subbundle (the corner reading) & M2 Def.~2.2 \\
$\tau,\ \hat\tau$ & flow-reversing involution of A2 ($\tau^{2}=$ fold on the cover); its lift & M2 \S2 \\
$S$, $v_{u}$, $v_{s}$, $\kappa=\log\varphi^{2}$ & quarter-turn; unstable/stable directions of $M_{0}$; Sol rate & M2 \S2 \\
$\omega=\psi+\bar\psi$ & the Coherent Field; $\hat\tau$-even, quaternion-valued, carries the imaginary direction & M2 Def.~2.3 \\
$\psi-\bar\psi$ & the odd sheet variable; $\hat\tau$-odd & M2 Def.~2.3 \\
$\QCD=\omega\lvert\psi-\bar\psi\rvert^{2}$ & the Quantum Coherence Dilaton, componentwise in $M$ & M2 Prop.~2.4 \\
$A$, $w=e^{A}$ & the A2 warp exponent and warp factor on the base & M2 Def.~2.5 \\
$k_{\omega}$ & decay rate of $A$ along $\theta_{\omega}$ & M2 \S4 \\
$F$ & curvature of the $\Spinc$ determinant line & M2 Def.~2.5 \\
$\kappa_{2}$ & two-dimensional gravitational coupling on the base & M2 \S4 \\
$T^{\QCD}_{MN}$ & stress tensor of $\QCD$, the module trace of the scalar form & M2 \S3 \\
$\zeta$ & coupling of $\QCD$: coefficient of $T^{\QCD}_{MN}$; equals the warp exponent per half-period & M2 \S\S3--4 \\
$G_{\mathrm{eff}},\ \Lambda_{\mathrm{eff}}$ & effective gravitational and vacuum couplings & M2 \S3 \\
\bottomrule
\end{tabular}
\end{center}

\paragraph{Introduced in M3.}
\begin{center}\small
\begin{tabular}{@{}p{0.29\textwidth}p{0.49\textwidth}p{0.15\textwidth}@{}}
\toprule
Symbol & Meaning & Where defined \\
\midrule
$\Sigma$-$0,\dots,\Sigma$-$8$, $\Sigma$-P, $\Sigma$-S, $\Sigma$-T & the twelve creation axioms of the earlier statement ($\Sigma$-3a a part of $\Sigma$-3) & Appendix~\ref{app:axioms} \\
$\mathbb{M}$-$1,\dots,\mathbb{M}$-$8$ & the eight mathematical axioms of the earlier statement & Appendix~\ref{app:axioms} \\
T, TT, A, N & the four classes: theorem of M1--M2; theorem target; the declared axiom; not an axiom & Definition~\ref{def:classes} \\
A5 & the fifth axiom (shadow) & Axiom~\ref{ax:A5}, Definition~\ref{def:A5} \\
$\rho$ & the reduction: integration over $\Keight$ & Definition~\ref{def:reduction} \\
$\dodd(n)$ & the number of odd divisors of $n$ & Section~\ref{sec:targets} \\
\bottomrule
\end{tabular}
\end{center}


\section{Machine-verified certificates}\label{app:lean}

Three certificates accompany the paper, each a single Lean~4 file on the core toolchain with no Mathlib dependency, no author axioms and no \texttt{sorry}. They build within the series' proof tree \texttt{MEFProof.Paper\_M3}; the build record is
\begin{lstlisting}
Built MEFProof.Paper_M3.M3_L1_lift   (496ms)
Built MEFProof.Paper_M3.M3_L2_count  (230ms)
Built MEFProof.Paper_M3.M3_L3_walls  (237ms)
\end{lstlisting}
Every certified statement is named at the point in the text where it is used. The certificates fix finite bookkeeping only. No classification is certified: each is argued in the text from a theorem of M1 or M2.

\texttt{M3\_L1\_lift.lean} (Section~\ref{sec:axiom}, 8 theorems): the fold on the grid $(\Z/4)^{2}$ is an involution with exactly four fixed points; it is two-to-one elsewhere; it is not one-to-one; the grid $(\Z/8)^{2}$ lies four-to-one over it compatibly with the fold; a count of assignments on four loops, not used in the text.

\texttt{M3\_L2\_count.lean} (Sections~\ref{sec:setting}, \ref{sec:result}, 6): the twenty axioms; the counts of the four classes and their sum; the two entries of class A as one axiom; the five sentences; the one scale.

\texttt{M3\_L3\_walls.lean} (Section~\ref{sec:theorems}, 7): the action of the fold, the drive and the exchange on the four corners; $S^{2}=-I$; the corner signs with product $-1$; the sum of the two handedness signs.

The three files follow verbatim.

\subsection*{\texttt{M3\_L1\_lift.lean} --- the fold in a finite model}
\lstinputlisting{../../MEFProof/Paper_M3/M3_L1_lift.lean}

\subsection*{\texttt{M3\_L2\_count.lean} --- the bookkeeping of the classification}
\lstinputlisting{../../MEFProof/Paper_M3/M3_L2_count.lean}

\subsection*{\texttt{M3\_L3\_walls.lean} --- corners, fold, drive and exchange}
\lstinputlisting{../../MEFProof/Paper_M3/M3_L3_walls.lean}

\begin{thebibliography}{99}
\bibitem{Anderson1972} P.~W. Anderson, More is different, \emph{Science} 177 (1972) 393.
\bibitem{Dirac1928} P.~A.~M. Dirac, The quantum theory of the electron, \emph{Proc. Roy. Soc. A} 117 (1928) 610.
\bibitem{Einstein1915} A. Einstein, Die Feldgleichungen der Gravitation, \emph{Sitzungsber. Preuss. Akad. Wiss.} (1915) 844.
\bibitem{Frobenius1877} G. Frobenius, \"Uber lineare Substitutionen und bilineare Formen, \emph{J. Reine Angew. Math.} 84 (1877) 1.
\bibitem{Gauss1827} C.~F. Gauss, Disquisitiones generales circa superficies curvas, \emph{Comm. soc. reg. sci. Gott. rec.} 6 (1827) 99.
\bibitem{Hamilton1844} W.~R. Hamilton, On a new species of imaginary quantities, connected with a theory of quaternions, \emph{Proc. Royal Irish Acad.} 2 (1844) 424.
\bibitem{Hilbert1900} D. Hilbert, Mathematische Probleme, \emph{G\"ottinger Nachrichten} (1900) 253.
\bibitem{Jacobi1829} C.~G.~J. Jacobi, \emph{Fundamenta Nova Theoriae Functionum Ellipticarum}, Borntr\"ager, K\"onigsberg, 1829.
\bibitem{Jordan1870} C. Jordan, \emph{Trait\'e des substitutions et des \'equations alg\'ebriques}, Gauthier-Villars, Paris, 1870.
\bibitem{Kaluza1921} T. Kaluza, Zum Unit\"atsproblem der Physik, \emph{Sitzungsber. Preuss. Akad. Wiss.} (1921) 966.
\bibitem{Luders1954} G. L\"uders, On the equivalence of invariance under time reversal and under particle-antiparticle conjugation for relativistic field theories, \emph{Danske Mat. Fys. Medd.} 28 (1954) No.~5.
\bibitem{Pauli1955} W. Pauli, Exclusion principle, Lorentz group and reflection of space-time and charge, in \emph{Niels Bohr and the Development of Physics}, Pergamon Press, London, 1955, p.~30.
\bibitem{Ramanujan1920} S. Ramanujan, Letter to G.~H. Hardy, 12 January 1920, in \emph{Collected Papers of Srinivasa Ramanujan}, Cambridge University Press, 1927.
\bibitem{Riemann1868} B. Riemann, \"Uber die Hypothesen welche der Geometrie zu Grunde liegen, \emph{Abh. K\"on. Ges. Wiss. G\"ottingen} 13 (1868) 133.
\bibitem{Weyl1918} H. Weyl, Gravitation und Elektrizit\"at, \emph{Sitzungsber. Preuss. Akad. Wiss.} (1918) 465.
\bibitem{Zwegers2002} S.~P. Zwegers, \emph{Mock Theta Functions}, PhD thesis, Utrecht, 2002.
\end{thebibliography}
\end{document}
